%% ============================================================
%%  resume.tex — Applied Mathematics Resume
%%  Fonts: Palatino  |  Icons: marvosym
%%  Compile: pdflatex resume.tex  (run twice)
%% ============================================================

\documentclass[10pt,letterpaper]{article}

%% ── Geometry ────────────────────────────────────────────────
\usepackage[top=0.65in,bottom=0.65in,left=0.70in,right=0.70in]{geometry}

%% ── Encoding & Font ─────────────────────────────────────────
\usepackage[T1]{fontenc}
\usepackage[utf8]{inputenc}
\usepackage{palatino}
\usepackage{microtype}

%% ── Math ────────────────────────────────────────────────────
\usepackage{amsmath,amssymb}

%% ── Colors ──────────────────────────────────────────────────
\usepackage[dvipsnames,svgnames]{xcolor}
\definecolor{NavyDeep}{HTML}{1B2A4A}
\definecolor{AccentTeal}{HTML}{2A6E7C}
\definecolor{RuleGray}{HTML}{C8CDD5}
\definecolor{SubtleGray}{HTML}{6B7280}

%% ── Hyperlinks ──────────────────────────────────────────────
\usepackage[hidelinks]{hyperref}
\hypersetup{colorlinks=true,urlcolor=AccentTeal,linkcolor=NavyDeep}

%% ── Lists ───────────────────────────────────────────────────
\usepackage{enumitem}
\setlist[itemize]{
  leftmargin=1.4em,itemsep=1.5pt,parsep=0pt,topsep=2pt,
  label={\color{AccentTeal}\textbullet}
}

%% ── Tables ──────────────────────────────────────────────────
\usepackage{booktabs,tabularx,array}

%% ── Icons ───────────────────────────────────────────────────
\usepackage{marvosym}
\usepackage{pifont}

%% ── Page-break control ──────────────────────────────────────
\usepackage{needspace}   % \needspace{Xpt} — keeps header+content together

%% ── Misc ────────────────────────────────────────────────────
\usepackage{parskip,setspace,calc,etoolbox,multicol}
\pagestyle{empty}

%% Suppress parskip glue between items we control manually
\setlength{\parskip}{0pt}

%% ============================================================
%%  CUSTOM COMMANDS
%% ============================================================

%% Section header: rule + title kept together with following content
%% \needspace{5\baselineskip} reserves ~5 lines; increase if sections
%% still break — header + first entry must both fit.
\newcommand{\cvsection}[1]{%
  \needspace{5\baselineskip}%
  \vspace{8pt}%
  {\color{NavyDeep}\rule{\linewidth}{0.8pt}\par}%
  \vspace{-2pt}%
  {\color{NavyDeep}\large\bfseries\MakeUppercase{#1}\par}%
  \vspace{3pt}%
}

%% Entry: \cventry{Title}{Org}{Date}{Details}
\newcommand{\cventry}[4]{%
  \noindent
  \begin{tabularx}{\linewidth}{@{}X r@{}}
    {\bfseries\color{NavyDeep}#1}\enspace{\color{SubtleGray}\itshape#2} &
    {\small\color{SubtleGray}#3}\\[-1pt]
    \multicolumn{2}{@{}l@{}}{\small\color{SubtleGray}#4}
  \end{tabularx}%
  \vspace{1pt}%
}

%% Publication: \cvpub{Authors}{Title}{Venue}{Year}{URL}
\newcommand{\cvpub}[5]{%
  \noindent{\small
    {\color{SubtleGray}#1.}\enspace
    ``{\itshape#2}.''
    \enspace{\bfseries#3},\enspace#4.%
    \ifstrempty{#5}{}{\enspace\href{#5}{{\color{AccentTeal}[link]}}}%
  }\par\vspace{3pt}%
}

%% Skill row
\newcommand{\cvskill}[2]{%
  \noindent{\small{\bfseries\color{NavyDeep}#1:\enspace}{\color{SubtleGray}#2}}\par\vspace{2pt}%
}

%% Course tag
\newcommand{\coursetag}[1]{%
  \mbox{\small\color{AccentTeal}\texttt{#1}}%
}

%% Inline math shortcut
\newcommand{\eq}[1]{$#1$}

%% ============================================================
%%  DOCUMENT
%% ============================================================
\begin{document}

%% ── Header ──────────────────────────────────────────────────
\begin{center}
  {\fontsize{22}{26}\selectfont\bfseries\color{NavyDeep} Your Full Name}\\[4pt]
  {\color{AccentTeal}\rule{3.5cm}{1.2pt}}\\[5pt]
  {\small\color{SubtleGray}
    \Email\enspace\href{mailto:you@university.edu}{you@university.edu}
    \quad\textbar\quad
    \Telefon\enspace +1\,(555)\,000-0000
    \quad\textbar\quad
    \href{https://yoursite.com}{yoursite.com}
    \quad\textbar\quad
    \href{https://github.com/yourhandle}{github.com/yourhandle}
    \quad\textbar\quad
    City, State
  }
\end{center}
\vspace{2pt}
\begin{center}
  {\small\color{NavyDeep}
    \textit{Research interests:}\enspace
    Numerical Analysis\enspace$\cdot$\enspace
    Scientific Computing\enspace$\cdot$\enspace
    Stochastic Processes\enspace$\cdot$\enspace
    Optimization\enspace$\cdot$\enspace
    Mathematical Modeling
  }
\end{center}

%% ── Education ───────────────────────────────────────────────
\cvsection{Education}

\cventry
  {Ph.D.\ in Applied Mathematics}
  {University of Somewhere}
  {Expected May 2026}
  {Advisor: Prof.\ Jane Doe\enspace|\enspace GPA: 3.95/4.0}
\begin{itemize}
  \item Dissertation: \textit{``Adaptive Multilevel Methods for Stochastic PDEs on Complex Domains''}
  \item Coursework: \coursetag{Numerical PDEs} \coursetag{Functional Analysis}
    \coursetag{Stochastic Calculus} \coursetag{Optimization} \coursetag{ML Theory}
\end{itemize}

\vspace{4pt}
\cventry
  {B.S.\ Mathematics (Honors), Minor Computer Science}
  {State University}
  {May 2020}
  {Summa Cum Laude\enspace|\enspace GPA: 3.98/4.0}
\begin{itemize}
  \item Senior thesis: \textit{``Finite Element Approximations to the Navier--Stokes Equations''}
\end{itemize}

%% ── Research Experience ─────────────────────────────────────
\cvsection{Research Experience}

\cventry
  {Graduate Research Assistant}
  {Computational Science Lab, University of Somewhere}
  {Aug 2020 – Present}
  {Advisor: Prof.\ Jane Doe}
\begin{itemize}
  \item Developed adaptive multilevel Monte Carlo (MLMC) for elliptic SPDEs achieving
        \eq{\mathcal{O}(h^2\varepsilon^{-2})} complexity — \eq{3\times} speedup over
        standard MLMC on benchmark problems.
  \item Implemented parallel C++/MPI solver on unstructured meshes; near-linear weak
        scaling to 512 cores on university HPC cluster.
  \item Applied sparse-grid collocation for UQ in subsurface flow models across
        20-dimensional parameter spaces.
\end{itemize}

\vspace{4pt}
\cventry
  {Research Intern — Mathematics of Machine Learning}
  {National Lab or Industry Partner}
  {Jun – Aug 2022}
  {Mentor: Dr.\ John Smith}
\begin{itemize}
  \item Proved sharp \eq{O(k^{-1/2})} SGD convergence under \eq{(L,\sigma)}-smooth
        assumptions; derived matching lower bounds.
  \item Benchmarked quasi-Newton methods for large-scale logistic regression,
        outperforming L-BFGS by 18\% on standard NLP datasets.
\end{itemize}

\vspace{4pt}
\cventry
  {REU Participant — Dynamical Systems}
  {Another University}
  {May – Jul 2019}
  {Mentor: Prof.\ Alice Brown}
\begin{itemize}
  \item Computer-assisted proof of Hopf bifurcation in a 3-species Lotka--Volterra
        model using interval arithmetic.
\end{itemize}

%% ── Publications ────────────────────────────────────────────
\cvsection{Publications \& Preprints}

\cvpub{Your Name, Jane Doe}
  {Adaptive Multilevel Monte Carlo for Stochastic Elliptic Equations}
  {SIAM Journal on Scientific Computing}{2024}
  {https://doi.org/10.xxxx/xxxxxx}

\cvpub{Your Name, John Smith}
  {Sharp Convergence Rates for Nonconvex SGD with Variance Reduction}
  {arXiv preprint arXiv:2401.00000}{2023}
  {https://arxiv.org/abs/2401.00000}

\cvpub{Your Name, Alice Brown}
  {Computer-Assisted Proof of Hopf Bifurcation in a Predator--Prey System}
  {Journal of Mathematical Biology}{2020}{}

%% ── Talks ───────────────────────────────────────────────────
\cvsection{Talks \& Presentations}

\noindent{\small
\begin{tabularx}{\linewidth}{@{}l X r@{}}
  \textbf{SIAM Annual Meeting}  & ``Adaptive MLMC for SPDEs'' (contributed talk)     & 2024\\[1pt]
  \textbf{Applied Math Seminar} & University of Somewhere (invited)                  & 2023\\[1pt]
  \textbf{ICIAM}                & ``Nonconvex SGD convergence rates'' (poster)        & 2023\\[1pt]
  \textbf{AMS Spring Sectional} & ``Hopf bifurcations via interval arithmetic'' (poster) & 2020\\
\end{tabularx}
}

%% ── Teaching ────────────────────────────────────────────────
\cvsection{Teaching Experience}

\cventry
  {Teaching Assistant — Numerical Analysis (MATH 5XX)}
  {University of Somewhere}
  {Fall 2021, Spring 2023}
  {Lead sections, write \& grade problem sets, hold office hours for 45 students}

\vspace{3pt}
\cventry
  {Instructor of Record — Calculus II (MATH 1XX)}
  {University of Somewhere}
  {Summer 2022}
  {Sole instructor; course evaluations 4.7/5.0}

%% ── Skills ──────────────────────────────────────────────────
\cvsection{Technical Skills}

\cvskill{Languages}{Python (NumPy, SciPy, JAX), C/C++, MATLAB, Julia, R, Bash}
\cvskill{HPC \& Parallel}{MPI, OpenMP, SLURM, GPU computing (CUDA basics)}
\cvskill{Math Software}{FEniCS, PETSc, deal.II, Firedrake, LAPACK/BLAS}
\cvskill{ML / Data}{PyTorch, scikit-learn, pandas, Jupyter}
\cvskill{Tools}{Git, \LaTeX{}, Linux, Docker, CMake}
\cvskill{Methods}{FEM, FDM, Spectral Methods, MLMC, Optimization, Bayesian Inference}

%% ── Honors ──────────────────────────────────────────────────
\cvsection{Honors, Fellowships \& Awards}

\noindent{\small
\begin{tabularx}{\linewidth}{@{}X r@{}}
  NSF Graduate Research Fellowship                           & 2021--2026\\[1pt]
  SIAM Student Paper Prize (honorable mention)               & 2023\\[1pt]
  Dean's Doctoral Excellence Award, University of Somewhere  & 2022\\[1pt]
  Phi Beta Kappa, State University                           & 2020\\[1pt]
  Barry Goldwater Scholarship                                & 2019\\
\end{tabularx}
}

%% ── Service ─────────────────────────────────────────────────
\cvsection{Service \& Professional Activities}

\noindent{\small\color{SubtleGray}
Referee: \textit{SIAM J.\ Numer.\ Anal.}, \textit{Math.\ Comp.},
\textit{J.\ Comput.\ Phys.}\quad$\cdot$\quad
Organizer: Graduate Applied Math Seminar, 2022--2024\quad$\cdot$\quad
Member: SIAM, AMS
}

\end{document}
